\hspace{3mm}

\centerline{\bf \Huge Теория по играм.}

\centerline{\bf \Huge Шутки и симметрии.}

\begin{flushright}
    \scriptsize{У нас нет времени для игр, мы должны воевать!\\ *удар грома*\\Shaman. Warcraft III: Reign of Chaos.}
\end{flushright}

\textbf{\Large Игры-шутки.}

Первый класс игр, о которых пойдет речь - игры-шутки. Это игры, исход которых не зависит от того, как играют соперники. Поэтому для решения такой игры не нужно указывать выигрышную стратегию. Достаточно лишь доказать, что выигрывает тот или иной игрок (независимо от того, как будет играть!). 

Приведем пример такой игры-шутки.

\textbf{Задача 1.} Двое по очереди ломают шоколадку 6 × 8. За ход разрешается сделать прямолинейный разлом любого из кусков вдоль углубления. Проигрывает тот, кто не    сможет сделать ход.

\textbf{Решение}. Основное соображение: после каждого хода количество кусков увеличивается ровно на 1. 

Сначала был один кусок. В конце игры, когда нельзя сделать ни одного хода, шоколадка разломана на маленькие дольки. А их 48! Таким образом, игра будет продолжаться ровно 47 ходов. Последний, 47-й ход (так же, как и все другие ходы с нечетными номерами) сделает первый игрок. Поэтому он в этой игре побеждает, причем независимо от того, как будет играть.

А теперь еще несколько игр-шуток:

\textbf{Задача 2.} Имеется три кучки камней: в первой - 10, во второй - 15, в третьей -20. За ход разрешается разбить любую кучку на две меньшие; проигрывает тот, кто не сможет сделать ход.

\textbf{Решение.}
Аналогично предыдущей задаче в конце игры мы получим 25 кучек камней, содержащих по одному камню. Всего будет сделано 23 хода (и это не зависит от того, как делают ходы игроки), следовательно, последний (нечетный) ход сделает первый игрок.


\textbf{Задача 3.} На доске написаны 10 единиц и 10 двоек. За ход разрешается стереть две любые цифры и, если они были одинаковыми, написать двойку, а если разными - единицу. Если последняя оставшаяся на доске цифра - единица, то выиграл первый игрок, если двойка - то второй.

\textbf{Решение.}
Чётность числа единиц на доске после каждого хода не меняется. Поскольку сначала единиц было чётное число, то после последнего хода на доске не может оставаться одна (нечётное число!) единица.

\textbf{\Large Симметрии}

\textbf{Задача 1.} Двое по очереди кладут пятаки на круглый стол, причем так, чтобы они не накладывались друг на друга. Проигрывает тот, кто не может сделать ход.

\textbf{Решение:} В этой игре выигрывает первый, независимо от размеров стола! Первым ходом он кладет пятак так, чтобы центры монеты и стола совпали. После этого на каждый ход второго игрока начинающий отвечает симметрично относительно центра стола. Отметим, что при такой стратегии после каждого хода первого игрока позиция симметрична. Поэтому если возможен очередной ход второго игрока, то возможен и симметричный ему ответный ход первого. Следовательно, он побеждает.

Попробуем, применяя симметричную стратегию, решить еще одну задачу.

\textbf{Задача 2.} Двое по очереди ставят слонов в клетки шахматной доски так, чтобы слоны не били друг друга. (Цвет слонов значения не имеет). Проигрывает тот, кто не может сделать ход.

\hspace{3mm}

\textit{Попытка решения:} поскольку шахматная доска симметрична относительно своего центра, то естественно попробовать симметричную стратегию. Но на этот раз (первым ходом нельзя поставить слона в центр доски) симметрию может поддерживать второй игрок. Казалось бы, по аналогии с предыдущей задачей, это и есть его выигрышная стратегия. Однако следуя ей, второму игроку не удастся сделать даже свой первый ход! Слон, только что поставленный первым игроком, может бить центрально-симметричное поле.

Этот пример показывает, что при доказательстве правильности симметричной стратегии нельзя забывать о следующем обстоятельстве: \textbf{очередному симметричному ходу может помешать ход. только что сделанный противником!} 

Ранее сделанные ходы, в силу симметрии позиции, помешать не могут. Что-бы решить игру при помощи симметричной стратегии, необходимо найти симметрию, при которой только что сделанный противником ход не препятствует осуществлению избранного плана.

Так, решение задачи легко провести, применяя не центральную, а осевую симметрию шахматной доски. За ось симметрии можно взять прямую, разделяющую четвертую и пятую горизонтали.

\textbf{\Large Виды симметрий в задачах на игры}

\hspace{3mm}

1)\textbf{ Геометрическая симметрия.}

Такую мы уже наблюдали в разобранных задачах. Сначала была симметрия относительно точки(\textit{центральная симметрия}), потом симметрия относительно прямой(\textit{осевая симметрия}).

2) \textbf{Повторюшка!}

Иногда бывает полезно посмотреть, что произойдет, если повторять ход за противником, как, например, в следующей задаче.

На столе лежат 2022 карточки с числами 1, 2, 3,... , 2022. Двое играющих берут по одной карточке по очереди. После того, как будут взяты все карточки, выигравшим считается тот, у кого больше последняя цифра суммы чисел на взятых карточках. Кто из играющих может всегда выигрывать, как бы ни играл противник, и как он должен при этом играть?

\textbf{Решение:} Ясно, что на результат влияют только последние цифры написанных на карточках чисел. Поэтому можно считать, что имеется по 202 карточки с цифрами 3, ..., 9, 0 и по 203 карточки с цифрами 1 и 2. Первый может сначала взять карточку с цифрой 2, а затем повторять ходы второго, пока это возможно. В момент, когда он не сможет повторить ход второго (после того, как второй возьмёт последнюю карточку с цифрой 1), первый берёт любую из оставшихся карточек, а далее снова повторяет ходы второго. И т. д. В результате в конце у обоих карточек каждого типа будет поровну,

\noindent
кроме карточек с цифрами 1 и 2: у первого будет на одну карточку с двойкой больше и на одну карточку с единицей меньше. Поскольку сумма  101·(0 + 1 + … + 9)  оканчивается пятеркой, сумма первого оканчивается семёркой, а второго – шестёркой.

3) \textbf{Дополни ход противника!} 

Так же бывает полезно посмотреть, а можно ли дополнить ход противника до чего-нибудь красивого или полезного. Следующая задача $-$ классический пример такой ситуации.

На столе лежат 40 бабуфохов. Герман и Эрдем могут взять из этой кучки либо 1 либо 2 бабуфоха. Проигрывает тот, кто не сможет ничего взять. Кто выиграет, если Герман ходит первым?

\textbf{Решение: } Выигрыает Герман. Первым ходом он берет ровно один бабуфох и тем самым оставляет 39(кратное трем) бабуфохов на столе. Теперь, если Эрдем берет 1 бабуфох, Герман берет 2 и наоборот. То есть дополняет ход Эрдема до 3 бабуфохов. А так как 39 делится на 3, то Герман заберет последние бабуфохи и выиграет. 

В следующий раз мы разберем, что такое выигрышные позиции, как их распознавать и как их находить с помощью метода анализа с конца.